\documentclass[12pt,a4paper]{article}

\usepackage[utf8]{inputenc}
\usepackage[T1]{fontenc}
\usepackage[spanish,es-nodecimaldot]{babel}
\usepackage{lmodern}
\usepackage{textcomp}
\usepackage{amsmath,amssymb}
\usepackage{booktabs}
\usepackage{array}
\usepackage{float}
\usepackage{xcolor}
\usepackage{microtype}
\usepackage{tikz}
\usepackage[hidelinks]{hyperref}
\usepackage[a4paper,margin=2.5cm]{geometry}

\usetikzlibrary{arrows.meta,babel,calc}

\setlength{\parindent}{0pt}
\setlength{\parskip}{0.65em}
\setcounter{tocdepth}{3}
\renewcommand{\arraystretch}{1.25}

\newcommand{\bitperiod}{T_b}

\begin{document}
\hypersetup{pageanchor=false}

\begin{titlepage}
\centering
\textbf{UNIVERSIDAD NACIONAL DE CÓRDOBA}\par
\vspace{0.5em}
\textbf{Facultad de Ciencias Exactas, Físicas y Naturales}\par
\vspace{0.5em}
\textbf{Ingeniería en Computación}\par
\vfill
{\LARGE\bfseries Trabajo Práctico N.\textsuperscript{o} 1\par}
\vspace{1.5em}
{\Large\bfseries Comunicaciones 101: repaso de fundamentos esenciales\par}
\vfill
\begin{tabular}{@{}ll@{}}
\textbf{Asignatura:} & Redes de Computadoras \\
\textbf{Grupo:} &  \\
\textbf{Alumnos:} &
  \begin{tabular}[t]{@{}ll@{}}
  Tomas Brigido & Kevin Cufre \\
  Rocco Delgado & Facundo Lugones \\
  Javier Siñuka & Agustin Tosolini
  \end{tabular} \\
\textbf{Fecha:} & Agosto de 2026
\end{tabular}
\vfill
\end{titlepage}

\hypersetup{pageanchor=true}
\pagenumbering{roman}
\tableofcontents
\clearpage
\pagenumbering{arabic}

\section{Fundamentos esenciales y análisis de una onda electromagnética}

\subsection{Fundamentos básicos}

\subsubsection{Ondas electromagnéticas}

Una onda electromagnética es una perturbación formada por campos eléctricos y
magnéticos variables que se propaga transportando energía. En el vacío, estas
ondas se desplazan a la velocidad de la luz, cuyo valor exacto es

\[
c=299\,792\,458\ \mathrm{m/s}.
\]

En comunicaciones, las ondas electromagnéticas permiten transportar
información tanto por medios guiados, como cables y fibra óptica, como por
medios no guiados, como el aire o el espacio libre. Sus parámetros principales
son la frecuencia, la longitud de onda, la amplitud, la fase y la velocidad de
propagación. La relación fundamental entre ellos es

\[
c=\lambda f,
\]

donde $c$ es la velocidad de propagación, $\lambda$ la longitud de onda y $f$
la frecuencia.

\subsubsection{Modulación y demodulación}

La \textbf{modulación} es el proceso por el cual una señal de información
modifica alguna característica de una señal portadora, generalmente senoidal,
para que pueda transmitirse eficientemente por un medio físico. Las
características modificadas suelen ser la amplitud, la frecuencia o la fase.

La \textbf{demodulación} es el proceso inverso: permite recuperar la
información original a partir de la señal modulada recibida. En comunicaciones
digitales, los bits se representan mediante cambios discretos de los parámetros
de la portadora.

\subsubsection{Señales de tiempo continuo}

Una señal de tiempo continuo está definida para todo instante dentro de un
intervalo y se representa matemáticamente como $x(t)$, donde $t$ puede tomar
cualquier valor real. Las señales analógicas son ejemplos típicos; una onda
senoidal utilizada como portadora de radiofrecuencia constituye un caso
clásico.

\subsubsection{Señales de tiempo discreto}

Una señal de tiempo discreto está definida únicamente en determinados
instantes, generalmente separados por un intervalo constante. Se representa
como $x[n]$, donde $n$ es un índice entero. Las señales digitales se asocian
frecuentemente con señales de tiempo discreto y amplitud discreta; por ejemplo,
una secuencia binaria formada por ceros y unos.

\subsection{Análisis del gráfico de la onda electromagnética}

En el gráfico del enunciado se observa una onda periódica cuya intensidad varía
con la distancia. A partir de las marcas del eje horizontal, la separación
entre dos puntos equivalentes consecutivos permite estimar

\[
\lambda=60\ \mathrm{mm}=0{,}060\ \mathrm{m}.
\]

La curva negra representa la oscilación de la onda. La línea roja de trazos
actúa como envolvente y muestra una disminución progresiva de la intensidad
máxima conforme avanza la propagación.

\subsection{Frecuencia y longitud de onda}

Considerando que la onda viaja exactamente a la velocidad de la luz,

\[
f=\frac{c}{\lambda}
=\frac{299\,792\,458\ \mathrm{m/s}}{0{,}060\ \mathrm{m}}
=4\,996\,540\,966{,}67\ \mathrm{Hz}.
\]

Por lo tanto, los resultados son

\[
\boxed{\lambda=60\ \mathrm{mm}},\qquad
\boxed{f\approx 5{,}0\ \mathrm{GHz}}.
\]

\subsection{Región del espectro electromagnético y banda}

Una onda de aproximadamente $5\ \mathrm{GHz}$ pertenece a la región de las
\textbf{microondas}. De acuerdo con la clasificación de la Unión Internacional
de Telecomunicaciones (ITU), el intervalo de $3$ a $30\ \mathrm{GHz}$
corresponde a la banda \textbf{SHF} (\emph{Super High Frequency}). También se
denomina región de ondas centimétricas, dado que sus longitudes de onda son del
orden de los centímetros.

\subsection{Dispositivos que operan en esta banda}

Un ejemplo representativo son los equipos de redes inalámbricas Wi-Fi que
operan en la banda de $5\ \mathrm{GHz}$, entre ellos:

\begin{itemize}
  \item routers Wi-Fi de doble banda;
  \item puntos de acceso inalámbricos;
  \item placas de red Wi-Fi;
  \item computadoras y teléfonos con conectividad IEEE 802.11 en $5\ \mathrm{GHz}$.
\end{itemize}

Como caso concreto puede considerarse un router Wi-Fi doméstico o empresarial
que conecta computadoras y teléfonos a una red local inalámbrica.

\subsection{Fenómeno representado por la línea roja de trazos}

La línea roja representa la \textbf{atenuación}: la disminución de la amplitud,
intensidad o potencia de una señal a medida que se propaga. En una transmisión
inalámbrica puede originarse por la distancia entre transmisor y receptor, la
dispersión de energía, la absorción causada por obstáculos, las reflexiones, la
difracción, la propagación multitrayectoria y las pérdidas propias del medio.

La amplitud de la onda queda progresivamente contenida por una envolvente
descendente, lo que ilustra que la señal llega cada vez con menor intensidad.

\subsection{Efecto sobre el dispositivo elegido}

La atenuación afecta a los dispositivos Wi-Fi de $5\ \mathrm{GHz}$. En la vida
cotidiana se observa cuando disminuye la señal al alejarse del router, se reduce
la velocidad al cambiar de habitación o paredes, muebles y estructuras
metálicas degradan el enlace. Habitualmente, una red de $5\ \mathrm{GHz}$ tiene
menor alcance que una de $2{,}4\ \mathrm{GHz}$, aunque pueda brindar mayor
velocidad a distancias cortas.

\subsection{Atenuación en distintos medios de transmisión}

\subsubsection{Telefonía celular}

La atenuación afecta a las transmisiones celulares. Las señales se propagan por
el espacio y atraviesan obstáculos como paredes, edificios, vegetación y
vehículos, por lo que la potencia recibida disminuye con la distancia y con las
pérdidas del entorno. Esto explica la menor cobertura dentro de edificios, en
subsuelos, ascensores, zonas rurales o lugares alejados de las antenas.

\subsubsection{Cable coaxial}

También existe atenuación en el cable coaxial. Aunque es un medio guiado, la
señal se debilita con la distancia debido a pérdidas eléctricas y dieléctricas,
que suelen crecer con la frecuencia y la longitud del tendido. En enlaces largos
pueden requerirse amplificadores, repetidores, cables de menor pérdida y
conectores adecuados.

\subsubsection{Fibra óptica}

La atenuación afecta asimismo a la fibra óptica. La potencia luminosa disminuye
por absorción, dispersión, curvaturas excesivas, imperfecciones del material y
pérdidas en empalmes y conectores. Aunque sus pérdidas son menores que las de
muchos medios eléctricos, la atenuación continúa siendo un parámetro central en
el diseño del enlace.

\section{Transmisión digital y comunicación tipo UART}

\subsection{Tipo y modo de transmisión representados}

En el diagrama del enunciado aparecen dos módulos conectados mediante líneas en
una única dirección. Una transporta los datos y la otra una señal periódica de
reloj. Por su direccionalidad, se trata de una transmisión \textbf{simplex}, ya
que la información circula solamente desde el emisor hacia el receptor. Por sus
características temporales, es \textbf{sincrónica}, porque el receptor emplea
una referencia de reloj para interpretar los datos. Los bits se envían en forma
secuencial, por lo que además es \textbf{serial}.

En síntesis, el sistema representa una transmisión
\textbf{serial, simplex y sincrónica}.

\subsection{Conveniencia para una transmisión rápida y bidireccional}

El esquema no es el más adecuado si se busca transmitir rápidamente y en ambos
sentidos, pues su principal limitación es que la comunicación es unidireccional.
Sería necesario agregar un canal en sentido contrario o compartir el medio con
un mecanismo de acceso apropiado.

Una alternativa sería un sistema \textbf{full-duplex}, si se desea transmitir y
recibir simultáneamente, o \textbf{half-duplex}, si se admiten ambos sentidos
pero no al mismo tiempo. En un enlace real de alta velocidad también deben
considerarse la sincronización, la codificación, la interferencia, el ruido, la
diafonía y las pérdidas del medio.

\subsection{Señal correspondiente a la cuarta letra del nombre del grupo}

Se adopta como nombre del grupo \textbf{Datosfera}; su cuarta letra es
\textbf{o}. El código ASCII correspondiente es

\[
111_{10}=\mathtt{6F}_{16}=01101111_2.
\]

Por lo tanto, el byte que debe representarse es $\mathbf{01101111}$. Siguiendo
el criterio del gráfico del enunciado, un nivel alto representa un uno y un
nivel bajo representa un cero. La señal ideal se muestra en la
Figura~\ref{fig:uart}.

\begin{figure}[H]
\centering
\begin{tikzpicture}[x=1.25cm,y=1.0cm]
  \draw[->] (-0.25,0) -- (8.45,0) node[right] {$t$};
  \draw[->] (0,-0.35) -- (0,1.65) node[above] {nivel};
  \draw[dashed,gray!70] (0,1) -- (8,1);
  \draw[dashed,gray!70] (0,0) -- (8,0);
  \node[left] at (0,1) {$1$};
  \node[left] at (0,0) {$0$};
  \foreach \x in {0,...,8}{
    \draw[gray!45] (\x,-0.10) -- (\x,1.35);
    \node[below] at (\x,-0.10) {$T_{\x}$};
  }
  \draw[very thick,blue!70!black]
    (0,0) -- (1,0) -- (1,1) -- (3,1) -- (3,0) --
    (4,0) -- (4,1) -- (8,1);
  \foreach \x/\b in {0/0,1/1,2/1,3/0,4/1,5/1,6/1,7/1}{
    \node[font=\bfseries] at (\x+0.5,1.48) {\b};
  }
\end{tikzpicture}
\caption{Representación ideal del byte $01101111$.}
\label{fig:uart}
\end{figure}

La correspondencia entre intervalos y niveles es la siguiente:

\begin{table}[H]
\centering
\begin{tabular}{ccc}
\toprule
Intervalo & Bit & Nivel ideal \\
\midrule
$T_0$ a $T_1$ & 0 & Bajo \\
$T_1$ a $T_2$ & 1 & Alto \\
$T_2$ a $T_3$ & 1 & Alto \\
$T_3$ a $T_4$ & 0 & Bajo \\
$T_4$ a $T_5$ & 1 & Alto \\
$T_5$ a $T_6$ & 1 & Alto \\
$T_6$ a $T_7$ & 1 & Alto \\
$T_7$ a $T_8$ & 1 & Alto \\
\bottomrule
\end{tabular}
\end{table}

En una UART real suele existir un bit de inicio, bits de datos, un posible bit
de paridad y uno o más bits de parada. Muchas implementaciones también envían
primero el bit menos significativo. No obstante, para mantener la coherencia
con el gráfico del enunciado, aquí se representa directamente el byte en el
orden indicado.

\subsection{Marcas temporales adecuadas para medir la señal}

No conviene medir exactamente en los instantes de transición entre bits, porque
allí el nivel de tensión cambia y puede presentar pendiente. Una medición en ese
momento podría interpretarse erróneamente como cero o uno.

La medición debe realizarse en una zona estable de cada intervalo,
preferentemente en su centro. Si la duración de cada bit es $T_b$, los instantes
recomendados son

\[
t_n=T_n+\frac{T_b}{2},\qquad n=0,1,\ldots,7.
\]

Así se maximiza el margen temporal respecto de ambos bordes y se reduce la
probabilidad de muestrear durante una transición.

\section{Transmisión inalámbrica y modulación}

\subsection{Limitaciones de transmitir una señal escalonada}

No es conveniente transmitir directamente por un medio inalámbrico una señal
escalonada o cuadrada de banda base. Una señal cuadrada ideal contiene muchos
componentes de frecuencia: para producir transiciones instantáneas necesita
armónicos de frecuencia arbitrariamente alta y, en el caso ideal, un ancho de
banda infinito. En la práctica esto no es realizable y puede causar
interferencia sobre otros sistemas.

Además, las antenas no irradian eficientemente señales de muy baja frecuencia o
con componente continua. Para lograr dimensiones prácticas y una radiación
eficiente se utiliza una portadora de radiofrecuencia. El canal inalámbrico,
además, es ruidoso y puede ser selectivo en frecuencia; el filtrado, la
propagación multitrayectoria y las interferencias deforman los flancos de la
señal y dificultan la decisión del receptor.

Finalmente, la regulación del espectro limita las frecuencias y los anchos de
banda utilizables. Por estos motivos se emplean técnicas de modulación que
adaptan los datos digitales a una portadora apta para propagarse por el medio.

\subsection{Técnica de modulación representada}

El gráfico del enunciado muestra una portadora senoidal cuya frecuencia cambia
según el valor del bit. La técnica es \textbf{FSK} (\emph{Frequency Shift
Keying}) y, dado que se emplean dos frecuencias para dos símbolos posibles, se
trata específicamente de \textbf{BFSK} (\emph{Binary Frequency Shift Keying}).
La amplitud se mantiene aproximadamente constante y la información se codifica
mediante el cambio entre dos frecuencias.

\subsection{Modulación de la secuencia propuesta}

La secuencia dada es

\[
\mathbf{01110110}.
\]

Se asigna una frecuencia baja $f_0$ al bit cero y una frecuencia alta $f_1$ al
bit uno. La representación conceptual resultante aparece en la
Figura~\ref{fig:bfsk}.

\begin{figure}[H]
\centering
\begin{tikzpicture}[x=1.25cm,y=1.15cm]
  \draw[->] (-0.15,0) -- (8.45,0) node[right] {$t$};
  \foreach \x in {0,...,8}{
    \draw[gray!45] (\x,-0.60) -- (\x,1.10);
    \node[below] at (\x,-0.60) {$T_{\x}$};
  }
  \foreach \x/\b in {0/0,1/1,2/1,3/1,4/0,5/1,6/1,7/0}{
    \node[font=\bfseries] at (\x+0.5,1.27) {\b};
  }
  \draw[very thick,blue!70!black,domain=0:1,samples=35]
    plot (\x,{0.43*sin(720*\x)});
  \draw[very thick,blue!70!black,domain=1:2,samples=50]
    plot (\x,{0.43*sin(1440*(\x-1))});
  \draw[very thick,blue!70!black,domain=2:3,samples=50]
    plot (\x,{0.43*sin(1440*(\x-2))});
  \draw[very thick,blue!70!black,domain=3:4,samples=50]
    plot (\x,{0.43*sin(1440*(\x-3))});
  \draw[very thick,blue!70!black,domain=4:5,samples=35]
    plot (\x,{0.43*sin(720*(\x-4))});
  \draw[very thick,blue!70!black,domain=5:6,samples=50]
    plot (\x,{0.43*sin(1440*(\x-5))});
  \draw[very thick,blue!70!black,domain=6:7,samples=50]
    plot (\x,{0.43*sin(1440*(\x-6))});
  \draw[very thick,blue!70!black,domain=7:8,samples=35]
    plot (\x,{0.43*sin(720*(\x-7))});
  \node[below=0.78cm] at (0.5,0) {$f_0$};
  \node[below=0.78cm] at (1.5,0) {$f_1$};
  \node[below=0.78cm] at (2.5,0) {$f_1$};
  \node[below=0.78cm] at (3.5,0) {$f_1$};
  \node[below=0.78cm] at (4.5,0) {$f_0$};
  \node[below=0.78cm] at (5.5,0) {$f_1$};
  \node[below=0.78cm] at (6.5,0) {$f_1$};
  \node[below=0.78cm] at (7.5,0) {$f_0$};
\end{tikzpicture}
\caption{Representación conceptual BFSK de la secuencia $01110110$.}
\label{fig:bfsk}
\end{figure}

\subsection{Otras técnicas de modulación}

Existen otras técnicas digitales basadas en modificar parámetros de una
portadora senoidal:

\begin{itemize}
  \item \textbf{ASK (\emph{Amplitude Shift Keying}):} modifica la amplitud de la
  portadora según el símbolo. Una variante simple es OOK (\emph{On-Off
  Keying}), donde la portadora está presente para un símbolo y ausente para el
  otro.

  \item \textbf{PSK (\emph{Phase Shift Keying}):} modifica la fase. BPSK, por
  ejemplo, utiliza dos fases diferentes para representar cero y uno.

  \item \textbf{FSK (\emph{Frequency Shift Keying}):} modifica la frecuencia.
  BFSK utiliza dos frecuencias, mientras que las variantes de orden superior
  emplean más frecuencias para representar más símbolos.

  \item \textbf{QAM (\emph{Quadrature Amplitude Modulation}):} combina cambios
  de amplitud y fase. Permite transmitir más bits por símbolo, aunque necesita
  mejores condiciones de señal para conservar una tasa de error baja.
\end{itemize}

\subsection{Tasa de error de bits y comparación de prestaciones}

La \textbf{tasa de error de bits}, o \textbf{BER} (\emph{Bit Error Rate}), es
la relación entre la cantidad de bits recibidos incorrectamente y la cantidad
total de bits transmitidos:

\[
\mathrm{BER}=\frac{\text{bits erróneos}}{\text{bits transmitidos}}.
\]

Por ejemplo, si se transmiten $1\,000\,000$ de bits y se reciben incorrectamente
$100$, se obtiene

\[
\mathrm{BER}=\frac{100}{1\,000\,000}=10^{-4}.
\]

Cuanto menor es el BER, mejor es el desempeño del enlace. En términos
generales, ASK es sensible al ruido de amplitud; FSK resulta robusta ante
variaciones de amplitud porque codifica la información en la frecuencia; y PSK,
en particular BPSK o QPSK con detección coherente, suele brindar mejores
prestaciones de BER para una misma relación entre energía por bit y densidad de
ruido en un canal ideal. QAM aumenta la cantidad de bits por símbolo, pero las
constelaciones de mayor orden exigen una relación señal-ruido más alta para
mantener el mismo BER.

Por ello, entre las técnicas básicas consideradas, \textbf{BPSK coherente suele
presentar mejores prestaciones de BER que ASK y FSK en un canal ideal con ruido
blanco gaussiano}. La elección real depende también del canal, el ancho de banda,
la complejidad del receptor y la potencia disponible.

\section{Simulacion en Packet Tracer}

\section{Conclusiones}

Los tres primeros puntos del trabajo relacionan los fundamentos físicos de las
comunicaciones con técnicas prácticas de transmisión digital. La onda analizada
tiene una longitud de onda aproximada de $60\ \mathrm{mm}$ y una frecuencia
cercana a $5\ \mathrm{GHz}$, por lo que pertenece a la región de microondas y a
la banda SHF según la clasificación de la ITU.

Toda transmisión real presenta atenuación, tanto en medios inalámbricos como en
medios guiados. En las señales digitales también es importante medir los niveles
en instantes estables, alejados de las transiciones.

Finalmente, una señal escalonada no es adecuada para su transmisión inalámbrica
directa por sus requisitos de ancho de banda, sus dificultades de radiación y
su sensibilidad a las distorsiones. Técnicas como FSK, ASK, PSK y QAM adaptan los
datos a portadoras electromagnéticas más apropiadas para el canal.

\begin{thebibliography}{9}

\bibitem{itu}
Unión Internacional de Telecomunicaciones,
\emph{Nomenclatura de las bandas de frecuencias y de las longitudes de onda
empleadas en telecomunicaciones}, Recomendación ITU-R V.431.

\bibitem{ieee80211}
IEEE,
\emph{IEEE 802.11: Wireless LAN Medium Access Control and Physical Layer
Specifications}.

\bibitem{catedra}
Bibliografía y apuntes de la cátedra de Redes de Computadoras sobre señales,
modulación, transmisión digital y medios físicos.

\end{thebibliography}

\end{document}
